\uuid{QKAB}
\titre{Ensembles}
\auteur{bodin ; croizat ; sacre}
\organisation{exo7}
\licence{CC BY-NC-SA 4.0}
\datecreate{2022-03-08}
\niveau{L1}
\module{Algèbre}
\chapitre{Logique, ensemble, raisonnement}
\sousChapitre{Ensemble}
\difficulte{3}
\theme{Ensembles}
\exo7qcmid{hilisit.logique.035}

\contenu{
\begin{question}
\qclassification{100.02}
Les fonctions $f$ et $g$ définies par les expressions suivantes sont-elles bijections réciproques l'une de l'autre ?
(On ne se préoccupera pas des ensembles de départ et d'arrivée.)
\begin{answers}
	\bad{$f(x) = \exp(2x)$ et $g(x) = \ln(\tfrac12x)$}
	\bad{$f(x) = \cos(x-1)$ et $g(x) = \sin(x+1)$}
	\good{$f(x) = \frac{1}{1+x}$ et $g(x) = \frac{1-x}{x}$}
	\bad{$f(x)= \sqrt{2x+1}$ et $g(x) = \frac12x^2-1$}
\end{answers}
\begin{explanations}
Si $f(x) = \exp(2x)$ et $g(x) = \ln(\tfrac12x)$ alors
$$(f \circ g)(x)
= e^{2 \ln(\tfrac12x) }
= e^{2 \ln(x) + 2\ln(\tfrac12) }
= e^{2 \ln(x))}\times e^{2\ln(\tfrac12) }
= x^2 \times (\tfrac12)^2 = \frac{x^2}{4}$$
et ne vaut donc pas $x$.

Si $f(x) = \cos(x-1)$ et $g(x) = \sin(x+1)$ alors $(f \circ g)(x) = \cos( \sin(x+1) -1)$ ne se simplifie pas du tout et n'a aucune chance d'être égal à $x$. Prendre par exemple $x=10$, alors $(f \circ g)(x)$ est une valeur renvoyée par un cosinus donc est compris entre $-1$ et $+1$ donc ne peut pas être égale à $10$.

Si $f(x) = \frac{1}{1+x}$ et $g(x) = \frac{1-x}{x}$
$$(f \circ g)(x)
= f\left(\frac{1-x}{x} \right)
=  \frac{1}{1+\frac{1-x}{x}}
= \frac{x}{x+(1-x)} = x.$$
On vérifie de même que $(g \circ f)(x) = x$. Donc $f$ et $g$ sont des bijections réciproques l'une de l'autre.

Si $f(x)= \sqrt{2x+1}$ et $g(x) = \frac12x^2-1$ alors
$$(f \circ g)(x)
= f\left(\tfrac12x^2-1  \right)
= \sqrt{ 2(\tfrac12x^2-1)+1 }
= \sqrt{ x^2-1}$$
qui n'est pas égale à $x$ (par exemple pour $x=10$, $(f\circ g)(10)= \sqrt{10^2-1}=\sqrt{99} \neq 10$).
\end{explanations}
\end{question}
}
